% Generated draft: requires mathematical and source review.
\documentclass[11pt]{article}
\usepackage[margin=25mm]{geometry}
\usepackage{fontspec}
\setmainfont{Latin Modern Roman}
\usepackage{amsmath,amssymb,mathtools,unicode-math}
\setmathfont{Latin Modern Math}
\usepackage{enumitem,needspace,xurl,hyperref}
\setlength{\parindent}{0pt}
\setlength{\parskip}{5pt}
\setlength{\emergencystretch}{3em}
\newcommand{\N}{\mathbb{N}}
\newcommand{\Z}{\mathbb{Z}}
\newcommand{\Q}{\mathbb{Q}}
\newcommand{\R}{\mathbb{R}}
\newcommand{\C}{\mathbb{C}}
\newcommand{\F}{\mathbb{F}}
\newcommand{\A}{\mathbb{A}}
\newcommand{\PP}{\mathbb{P}}
\newcommand{\E}{\mathbb{E}}
\newcommand{\eps}{\varepsilon}
\newcommand{\dd}{\,\mathrm{d}}
\begin{document}
\section*{UnsolvedMath ID: 6200032}
\textbf{Source problem code:} catalogue entry
\subsection{Definitions and mathematical statement}
Let $G$ be finitely generated, choose a finite generating set, and let $\Gamma$ be its Cayley graph with the path metric. A group is virtually a surface group if it has a finite-index subgroup isomorphic to the fundamental group of a surface in the intended class. The source does not specify whether surfaces must be closed, orientable, or of negative Euler characteristic, or whether the torus is included.

The supplied implication is
\[
 \Gamma\text{ admits a sequence of separating quasicircles}
 \quad\Longrightarrow\quad
 G\text{ is virtually a surface group}.
\]
Here ``quasicircle'' and ``separate'' are not defined. One metric possibility would use images of circles with controlled quasi-isometry constants, but the source does not prescribe that model, require uniform constants for the sequence, or explain whether the curves are embedded cycles. Topological disconnection by deleting a cycle and coarse separation into deep or unbounded components are different conditions.

The sequence also has no stated growth, scale, or exhaustiveness requirement: the record does not say whether diameters tend to infinity, whether repeated curves are excluded, or how far-separated regions must be. These quantitative data are essential to the intended rigidity question and cannot be inferred from the word ``sequence.'' The precise hypothesis is therefore retained as an unresolved source dependency. No one-endedness, hyperbolicity, finite-presentation, or other extra group assumption is added.
\subsection{Short English statement}
Does a finitely generated group whose Cayley graph has the intended sequence of separating quasicircles have a finite-index surface subgroup? A precise formulation needs the quasicircle control, separation depth, sequence-scale condition, and permitted surface class.

The supplied summary reports no resolution and connects the question to Papasoglu's quasiline/quasisphere separation work and surface-group rigidity, without a precise theorem or reference. The OPEN-TRIAGE label and historical 2005 problem are retained with current status explicitly marked for review.
\subsection{Sources}
\begin{itemize}
\item Local UnsolvedMath registry, permanent ID 6200032, source code catalogue entry, ``Boundaries of Groups and Kleinian Groups --- Problem 32,'' proposed by Panos Papasoglu, 2005. Dataset curation and metadata: Ulam AI and UnsolvedMath Contributors, CC BY 4.0.
\item Kapovich, \emph{Problems on Boundaries of Groups and Kleinian Groups} (2005), Problem 32, PDF page 9, \url{https://www.math.ucdavis.edu/~kapovich/EPR/problems.pdf}; normalized pdf-text extraction, accessed 2026-07-29.
\item The supplied summary mentions Papasoglu's separation work without complete bibliographic details or a checked implication. Source-specific provenance marks current status and statement-reuse terms NEEDS\_REVIEW for the public author-hosted PDF, separately from the dataset licence.
\end{itemize}
\end{document}
